% P10-4c-Nachtrag, 15.09.2026: Auswertungsregel für technische Handlungen
% und die bedingte Aussagekraft fehlender Ereignisse expliziert.
% P10-4c, 14.09.2026: Ledger-Herleitung und Nachweisanschlüsse präzisiert.
% M15, 11.09.2026: Ziel und offener Lösungsweg, Anforderungsherkunft
% sowie Belegarten getrennt erklärt; Instrumentierung kurz definiert.
% M03 / B-17, 10.09.2026: ursprüngliches und erweitertes Nutzungsziel
% unterschieden; keine neue Datierung des Betreuerimpulses.
% ============================================================
% §3.2 Iteratives Vorgehen und Evidenzgrundlage — v03 (Claude)
% v03 (01.09., Kollegen-Feedback nach Kap-4-Konkretisierung):
%   Punkt 1 (Instrumentierungs-Herkunft) war bereits umgesetzt
%   (Beobachtungsschicht-Absatz) — ergänzt nur „über OpenTelemetry".
%   Punkt 2 ÜBERNOMMEN: „Systemverhalten" pauschal→Kap.7 war zu weit
%   (widersprach dem Test-Satz davor); jetzt getrennt: VERGLEICHENDE/
%   Wirkungs-Aussagen → Kap.-7-Evaluation · technische Eigenschaften →
%   Verifikation durch Tests/Prüffälle (= 4.8-Vokabular, drei Ebenen
%   Beobachtung/Verifikation/Evaluation).
%
% v02 (Review-Runde 1 — kritisch geprüft):
%   GEGENSTANDSLOS: ① "alte Verweise in 3.1" — Kollege reviewte die
%     Fassung VOR der v03-Umstellung; aktuelle Labels sind korrekt
%     (per grep verifiziert). Sein Auftakt-/Absatz-3-Ersatz hätte
%     zudem die vom Autor beauftragte POSITIVE DREHUNG rückgängig
%     gemacht → nicht übernommen.
%   ÜBERNOMMEN (alle epistemisch berechtigt):
%   ② "repräsentative" → "projektnahe Eingaben" (Repräsentativität
%     wäre eine unbelegte Stichproben-Behauptung; real: echte/
%     projektnahe Transkripte — Abweichung vom Blaupause-Wortlaut
%     "repräsentativen" ist BEWUSST, Blaupause-v02 §7 ggf. nachziehen).
%   ③ "vielfach" → "wiederholt durchlaufen; überarbeitet, stabilisiert
%     oder verworfen" (keine Pseudo-Quantität).
%   ④ Dritter Herkunftstyp als "abgeleitete Realisierungsanforderungen"
%     + Anti-Zirkel-Satz — sein bester Punkt; liegt näher am
%     Blaupause-Wortlaut ("konsistente Realisierung, Governance,
%     Betrieb") als meine v01-Formulierung.
%   ⑤⑥⑦ Schlussabsatz konsolidiert: "dokumentieren" statt "belegen";
%     Tests = Evidenz unter Prüfbedingungen; und der wichtigste Fang:
%     auch Kapitel 7 liefert KEINE "allgemeinen" Wirksamkeitsaussagen,
%     sondern begrenzte Aussagen unter definierten Bedingungen (meine
%     v01-Formulierung hatte das implizit versprochen).
%     ("Traceability" aus seinem Vorschlag → "Nachvollziehbarkeit",
%     deutscher Begriffshaushalt von Kap. 2.)
%   ZURÜCKGEWIESEN: Zweifel am "Fehler- und Beobachtungslog" — das Log
%     EXISTIERT real und systematisch (fortlaufendes Reibungslog mit
%     Beleg-Lauf, Befund und Fix-Status je Eintrag; Projektpflichtregel
%     seit Beginn der E2E-Phase). Keine "verstreuten Notizen".
%   MERKER (berechtigt): Dateikopf-Kommentare vor finaler Abgabe aus
%     allen .tex entfernen (reine Arbeitsdokumentation).
%
% Grundlage: Kapitel-3-Blaupause v02 §7 (Absatzplan 1–6). KEINE neuen
% Methodenquellen (Blaupause: "Das eigene Vorgehen wird durch
% Projektdokumentation und Artefakte gestützt"); Hevner/Peffers wirken
% aus 3.1 nach.
%
% Realitäts-Abgleich (Claude, gegen die Projektgeschichte):
%   · M03/B-17, Autorenpräzisierung 10.09.2026: ursprüngliches Ziel
%     MAF + KI-Agenten + Kick-off-Transkript -> GitHub-Issue;
%     Rückkopplung/Wiederverwendung = Zielerweiterung. Die Architektur
%     war damit nicht vorgegeben. Der Betreuerimpuls ist nicht durch
%     zeitnahe Aufzeichnungen unabhängig datiert; keine Datierung
%     oder wörtliches Gesprächszitat aus Implementierungsdaten ableiten.
%   · Der 6-Schritte-Zyklus = der tatsächliche Arbeitsmodus (Stufe
%     bauen → mit echten Transkripten laufen lassen → beobachten →
%     Konsequenz → umbauen/stabilisieren).
%   · "fortlaufend geführtes Fehler- und Beobachtungslog" = das
%     Reibungslog (systemneutral formuliert; vom Autor-Vorschlag
%     31.08. gedeckt — beschreibt das markanteste formative
%     Instrument präziser als nur "Notizen").
%   · Alle genannten Evidenzformen existieren real (Versionshistorie,
%     Runs/Logs/Tool-Aufrufe, Artefakte, Tests, Notizen).
%
% Selbstcheck gegen Blaupause-§7-Abnahme:
%   Zielkorridor/Lösungsraum getrennt ✓ · Zyklus abstrakt, keine
%   Pivots/Agenten/Fehler ✓ · 3 Herkunftstypen ✓ · Auswahlregel +
%   Schutzregeln ✓ · Evidenzgrenzen als Prosa statt Tabelle (Blaupause
%   erlaubt "drei bis vier Sätze") ✓ · verbindlicher Grenzsatz wörtlich
%   getragen ✓ · Übergang zu Kap. 4 ✓ · ~440 Wörter ≈ 1,2 Seiten
%   (Budget 1,1–1,4) ✓ · systemneutral ✓ · keine Mess-Verben ✓.
% ============================================================

\section{Iteratives Vorgehen und Evidenzgrundlage}
\label{sec:iteratives-vorgehen-evidenz}

Ausgangspunkt war das Ziel, mit dem Microsoft Agent Framework
KI-Agenten einzusetzen, die ein Gesprächstranskript verarbeiten und
daraus ein GitHub-Issue ableiten. Später kam das Ziel hinzu, die
eingesetzten agentischen Fähigkeiten nach der Issue-Erzeugung erneut
zu nutzen und Rückmeldungen in einen Verarbeitungskreislauf
einzubeziehen. Damit rückte die Verarbeitung weiterer
Projektkommunikation in den Blick. Offen blieben die Rollen und
Grenzen der Agenten, ihre Orchestrierung, die Zwischenrepräsentationen,
das Zustandsmodell und die Aufteilung zwischen agentischen und
deterministischen Komponenten. Exploriert wurde der Lösungsweg zu
diesem im Verlauf erweiterten Nutzungsziel innerhalb der frühen
SDLC-Phasen. Die Implementierung der fachlichen Zielanwendung blieb
außerhalb des Untersuchungsgegenstands.

Die Entwicklung folgte einem wiederkehrenden Zyklus: Eine Teillösung
wurde entworfen oder verändert, technisch umgesetzt und mit
projektnahen Eingaben ausgeführt. Die erzeugten Artefakte und das
Laufzeitverhalten wurden beobachtet. Aus der Interpretation dieser
Beobachtungen wurden Konsequenzen für den Entwurf abgeleitet und die
Teillösung anschließend überarbeitet, stabilisiert oder verworfen.
Kapitel~\ref{chap:explorative-problemklaerung} behandelt die ausgewählten
Beobachtungen und die daraus hervorgegangenen Entscheidungen.

Die dort hergeleiteten Designanforderungen haben drei unterschiedliche
Ursprünge. \emph{Zielgetriebene Anforderungen} konkretisieren das
ursprüngliche oder erweiterte Nutzungsszenario und die Systemgrenze.
\emph{Anforderungen aus formativen Beobachtungen} gehen auf Probleme
oder Grenzen früher Entwicklungsstände zurück.
\emph{Abgeleitete Realisierungsanforderungen} ergeben sich aus der
technischen Umsetzung, dem Betrieb oder der Absicherung bereits
begründeter Architekturentscheidungen. Dieser dritte Typ konkretisiert
eine Entscheidung, liefert aber keine unabhängige Begründung für sie.
Die Unterscheidung hält gesetzte Ziele, beobachtete Probleme und
spätere Entscheidungen auseinander. So werden weder alle Anforderungen
nachträglich aus Fehlschlägen hergeleitet noch spätere Entscheidungen
als von Beginn an feststehend dargestellt.

Die Rekonstruktion stützt sich auf die im Projekt entstandenen Belege:
Quellcode und Versionshistorie, identifizierbare Entwicklungsstände,
Laufzeitereignisse und protokollierte Werkzeugaufrufe, Zwischen- und
Ergebnisartefakte, gezielte explorative Vorstudien mit dokumentierten
Grenzen sowie automatisierte Tests und zeitnahe Entwicklungsnotizen.
Zu Letzteren gehört ein fortlaufend geführtes Fehler- und
Beobachtungslog.

Die Laufzeitbelege entstanden durch eine schrittweise aufgebaute
Instrumentierung, also die technische Erfassung des
Ausführungsverhaltens. Eine anwendungsseitige Beobachtungsschicht
protokollierte laufbezogene Metadaten sowie Modell-, Workflow- und
Werkzeugereignisse, ergänzt um Telemetriesignale über OpenTelemetry.
Damit ließen sich sprachliche Modellausgaben, angeforderte
Werkzeugaktionen und beobachtbare Zustandswirkungen getrennt
betrachten. Der Umfang dieser Erfassung wurde während der Entwicklung
erweitert und war nicht in allen historischen Ständen gleich.
Aussagen aus Laufprotokollen werden deshalb auf die tatsächlich
erfassten Ausführungsschritte begrenzt. Die technische Realisierung
beschreibt Abschnitt~\ref{sec:realisierung-observability}.

Für Aussagen über ausgeführte Handlungen haben die zugehörigen
Werkzeug- und Systemereignisse Vorrang vor widersprechenden
Selbstauskünften des Modells; behauptete Zustandsänderungen werden
zusätzlich an den Ergebnisartefakten geprüft. Aus einem fehlenden
Ereignis wird nur dann auf eine unterbliebene Handlung geschlossen,
wenn der betreffende Ausführungspfad nachweislich instrumentiert
und der relevante Protokollabschnitt vollständig erhalten ist.
Andernfalls wird die fehlende Aufzeichnung als Nachweislücke behandelt.

Die Auswertung umfasst ausgewählte Entwicklungsbefunde. Aufgenommen
wird eine Beobachtung, wenn sie eine zentrale Designanforderung
motivierte, eine tragende Architekturentscheidung beeinflusste oder
eine für die Forschungsfragen relevante Grenze sichtbar machte.
Fehlende Aufzeichnungen werden nicht durch plausible Annahmen ersetzt
und nicht betrachtete Alternativen nicht als verworfen dargestellt.
Relevante Gegenbeobachtungen werden bei der jeweiligen Aussagegrenze
berücksichtigt.

Die Belegarten erfüllen unterschiedliche Aufgaben. Entwicklungsunterlagen
und formative Beobachtungen begründen die Rekonstruktion von
Gestaltungsentscheidungen. Einzelne Beobachtungen erlauben weder
Häufigkeitsangaben noch einen Nachweis der Wirksamkeit des finalen
Systems. Quellcode und Konfigurationen dokumentieren implementierte
Regeln; Tests und ausgeführte Prüffälle können technische Eigenschaften
unter den jeweils angegebenen Bedingungen verifizieren. Daraus folgt
noch keine Aussage über die fachliche Qualität der erzeugten Artefakte.

Formative Vergleiche können innerhalb ihrer jeweiligen
Versuchsbedingungen Verbesserungen und Grenzen sichtbar machen. Sie
begründen die weitere Gestaltung, ohne damit die Leistungsfähigkeit des
späteren Systems nachzuweisen. Dessen ausgewählte Eigenschaften werden
in Kapitel~\ref{chap:evaluation} bewertet; die dortigen Aussagen zu
Qualität, Nachvollziehbarkeit, Aufwand und Mechanismenwirkung gelten
für die ausgewiesenen Lauf-, Bewertungs- und Untersuchungsstände.
Die Projektbelege erlauben keine allgemeinen Aussagen über LLM-basierte
Agenten, das Requirements Engineering oder das verwendete Framework.
